% SOURCE_AUTHORITY: sga5_fr_workpass.tex, lines 6998--7031 (LF-normalized unit).
% SOURCE_SHA256: 791F4EFFC5E02832D5D77ED03518C8156D6F07E4C8238B03545DB93D883FBB28
\subsubsection*{2.4.3}
Vimos (2.4.1) que o subconjunto $AR$ de $\Fl(P)$ permite um cálculo de frações pela direita. Isso nos permite (2.3.3) considerar a categoria
\[
P'_{AR}=P(AR^{-1})
\]
que, lembremos, tem os mesmos objetos que $P$. Sejam $X$ e $Y$ dois objetos de $P$, e denotemos por $\Hom'_{AR}(X,Y)$ o conjunto dos isomorfismos de $X$ em $Y$ em $P'_{AR}$. Um elemento de $\Hom_{AR}(X,Y)$ define de maneira evidente um elemento de $\Hom'_{AR}(X,Y)$. Mais precisamente, define-se assim um funtor
\[
m:P_{AR}\longrightarrow P'_{AR}.
\]

\begin{statement}{Proposição 2.4.3}
O funtor $m$ é um isomorfismo de categorias.
\end{statement}

\begin{prooftext}[Demonstração]
A questão é ver que, para cada par $(X,Y)$ de objetos de $P$, a correspondente aplicação
\[
m_{XY}:\Hom_{AR}(X,Y)\longrightarrow\Hom'_{AR}(X,Y)
\]
é uma bijeção. Evidentemente, é uma sobrejeção. Para ver a injetividade, reduzimo-nos imediatamente a verificar que, para todo diagrama comutativo do tipo que aparece a seguir, existe um inteiro $t$ maior ou igual a $r$ e $s$ tal que a composição de $\rho$ e do morfismo canônico de $X[t]$ em $X[s]$ seja o morfismo canônico de $X[t]$ em $X[r]$:
\[
\begin{tikzcd}[column sep=large,row sep=large]
X[r] \arrow[d,"{v_{rX}}"'] & X[s] \arrow[l,"\rho"'] \arrow[dl,"{v_{sX}}"]\\
X & {}
\end{tikzcd}
\]
Ora, isso se deduz da comutatividade do diagrama que aparece a seguir, no qual se designaram pela letra $c$ todos os morfismos canônicos que intervêm:
\[
\begin{tikzcd}[column sep=large,row sep=large]
X[r+s] \arrow[d,"c"'] & X[2s] \arrow[l,"\rho(s)"] \arrow[d,"c"] \arrow[dl,"c" description]\\
X[r] & X[s] \arrow[l,"\rho"]
\end{tikzcd}
\]
\end{prooftext}
